
To prove the consistency of the IJ covariance,
we assume the following conditions on the model
and data generating process.

%%%%%%%%%%%%%%%%%%%%%%%%%%%%%%%%%%%%%%%%%
%%%%%%%%%%%%%%%%%%%%%%%%%%%%%%%%%%%%%%%%%
\begin{assu}\assulabel{bayes_clt}

%For the model and data generating process,
For $\delta > 0$, denote the $\delta$--neighborhood of
$\thetatrue$ as $\thetaball{\delta} :={}
\left\{\theta: \norm{\theta - \thetatrue}_2 \le \delta \right\}$.
Assume there exists a $\delta > 0$ such that:

%
\begin{enumerate}
%
\item \itemlabel{finite_dim} The dimension of $\theta$ stays fixed as $N$ changes.
%
\item \itemlabel{prior_smooth} The prior density $\prior(\theta)$ is non-zero
and bounded for all $\theta \in \thetadom$, and four times continuously
differentiable in $\thetaball{\delta}$.
%
\item \itemlabel{prior_proper}
The prior is proper.%
%
\item \itemlabel{loglik_smooth} The functions $\lik(\theta)$ and $\theta \mapsto
\ell(\xn \vert \theta)$ are four times continuously differentiable in
$\thetaball{\delta}$, and the interchange of partial differentiation
with respect to $\theta$ and expectations with respect to $\xfdist$ is
justified.
%
\item \itemlabel{loglik_ulln} There exists a function $M(\xn)$ such that
%$\deltaulln > 0$ such that
%
\begin{align*}
%
\sup_{\theta \in \thetaball{\delta}}
\norm{
\ellgrad{4}(\xn \vert \theta)
}_2 \le M(\xn)
\quad\textrm{and}\quad
\expect{\fdist(\xn)}{M(\xn)^2} < \infty.
%
\end{align*}
%
\item \itemlabel{strict_opt}
The log likelihood has a strict optimum in the following sense:
    %
    \begin{itemize}
    \item
    A unique maximum
    $\thetatrue = \mathrm{argmax}_{\theta \in \thetadom} \lik(\theta)$
    exists.
    %
    \item The matrix $\info$ is positive definite.
    %
    % \item For any $\delta > 0$, with probability approaching one as $N
    \item There exists an $\varepsilon > 0$ such that, with probability
    approaching one as $N \rightarrow \infty$,
    %
    %
    % This formulation is not obviously enough to guarantee that
    % \thetahat \in \thetaball{\delta} ... but it's required
    % for R3, since if the prior goes to zero, the likelihood
    % can blow up outside the delta ball and still have the MAP
    % be well-behaved.
    \begin{align*}
    \sup_{\theta: \thetadom \setminus \thetaball{\delta}}
    \left( \meann \ell(\x_n \vert \thetatrue) -
           \meann \ell(\x_n \vert \theta)
    \right) > \varepsilon.
    \end{align*}
    %
    % $\sup_{\theta: \thetadom \setminus \thetaball{\delta}}
    % \left( \likhat(\thetatrue) - \likhat(\theta)
    % \right) > \varepsilon$.
    %
    \end{itemize}
%
\end{enumerate}

\end{assu}
%%%%%%%%%%%%%%%%%%%%

% Under \assuref{bayes_clt} \itemref{loglik_ulln}, a ULLN will apply to
% $\ellgrad{k}(\x_n \vert \theta)$ for $k \in \{0, 1, 2, 3\}$ (see a proof in
% \lemref{ulln} of \appref{lemmas}).

Under \assuref{bayes_clt}, a BCLT will apply in total variation distance by
classical results (\citet{van:2000:asymptotic} Chapter 10), but we require an
expansion of posterior moments to analyze \eqref{ij_heuristic}.  Series
expansions of posterior moments are studied by
\citet{kass:1990:posteriorexpansions} (under stricter differentiability
conditions than ours), but with no control over the data dependence in the error
when the expectation depends on the data.

We now state conditions on functions $\phi(\theta)$ and $\phi(\theta, \x_n)$
sufficient to form a controllable series expansion of their posterior
expectations.

\begin{defn}\deflabel{bclt_okay}
%
A vector--valued function $\phi(\theta, \xn)$, possibly depending on a single
datapoint, is ``$K$-th order BCLT-okay'' if
%
\begin{enumerate}
%
\item\itemlabel{bclt_okay_as}
    $\theta \mapsto \phi(\theta, \xn)$
    is $K$ times continuously differentiable $\fdist$--almost surely.
%
\item\itemlabel{prior_op1}
    $\meann \expect{\prior(\theta)}{\norm{\phi(\theta, \x_n)}^2_2} = O_p(1)$.
%
\item\itemlabel{grad_op1}
For some $\delta > 0$, and for each $k = 1, \ldots, K$,
$\sup_{\theta \in \thetaball{\delta}}
    \meann \norm{\phigrad{k}(\theta, \x_n)}^2_2 = O_p(1)$.
%
\end{enumerate}
%
When the function $\phi(\theta, \x_n) = \phi(\theta)$ does not depend on the
data (e.g., $\phi(\theta, \x_n) = \g(\theta)$), then \itemref{grad_op1}
is vacuous, and \itemref{prior_op1} simply becomes
$\expect{\prior(\theta)}{\norm{\phi(\theta)}^2_2} < \infty$.
%
\end{defn}


%%%%%%%%%%%%%%%%%%%%%%%%%%%%%%%%%%%%%%%%%
%%%%%%%%%%%%%%%%%%%%%%%%%%%%%%%%%%%%%%%%%
%%%%%%%%%%%%%%%%%%%%%%%%%%%%%%%%%%%%%%%%%
%
\begin{thm}\thmlabel{bayes_clt_main}(See proof in \appref{bayes_clt_proof}.)
%
Let \assuref{bayes_clt} hold, and assume that $\phi(\theta, \x_n)$ is
third-order BCLT-okay.  Let $\normfourth$ denote the
$\thetadim\times\thetadim\times\thetadim\times\thetadim$ array of fourth moments
of the Gaussian distribution $\normdist(0, \infohat^{-1})$. With
$\fdist$-probability approaching one, for each $n$,
%
\begin{align}
%
\expect{\post}{\phi(\theta, \x_n)} - \phi(\thetahat, \x_n)
={}&
N^{-1} \left(
    \frac{1}{2} \phigrad{2}(\thetahat, \x_n) \infohat^{-1} +
    \frac{1}{6} \phigrad{1}(\thetahat, \x_n) \likhatk{3}(\thetahat) \normfourth
\right) +
\nonumber\\&
\ordlog{N^{-2}} 1_{\phi} \resid{\phi}(\x_n),
\textrm{ where }\meann \resid{\phi}(\x_n)^2 = O_p(1).
\eqlabel{bclt_expansion}
%
\end{align}
%
% where $\meann \resid{\phi}(\x_n)^2 = O_p(1)$.

In \eqref{bclt_expansion}, $\resid{\phi}(\x_n)$ is a scalar,
$1_{\phi}$ is a vector of ones with the same dimension as $\phi(\theta, \x_n)$, and
derivatives with
respect to $\theta$ are summed against the indices of the Gaussian
arrays.\footnote{Although a slight abuse of notation, these simple conventions
avoid a tedious thicket of indicial sums. Concretely, we mean that
%
\begin{align*}
% $\expect{\normdist(\tau \vert 0, \infohat^{-1})}{\tau_a \tau_b}$
\phigrad{2}(\thetahat, \x_n) \infohat^{-1} ={}&
    \sum_{a,b=1}^{\thetadim}
    \fracat{\partial^2 \phi(\theta, \x_n)}
           {\partial\theta_a \partial\theta_b}{\thetahat}
   (\infohat^{-1})_{ab}
\\
\phigrad{1}(\thetahat, \x_n) \likhatk{3}(\thetahat) \normfourth
={}&
\sum_{a,b,c,d=1}^{\thetadim}
\fracat{\partial \phi(\theta, \x_n)}{\partial\theta_a }{\thetahat}
\fracat{\partial^3 \lik(\theta)}
       {\partial\theta_b \partial\theta_c \partial \theta_d }{\thetahat}
       \normfourth_{abcd}
%
\end{align*}
%
each of which is a vector expression of the same length as $\phi(\theta,
\x_n)$.
}
%
\end{thm}
%
% \begin{proofsketch}{\appref{bayes_clt_proof}}
% %
% Our proof of \thmref{bayes_clt_main} follows other classical proofs of the BCLT
% (e.g. \cite[Theorem 1.4.2]{ghosh:2003:bayesian}) by dividing $\thetadom$ into
% three regions, only one of which --- a shrinking ball centered at $\thetahat$ ---
% contributes leading order terms to the posterior expectation. Within this
% dominant shrinking ball, we can Taylor expand the posterior integrals and
% collect terms to give the desired result.
% %
% Our distinct contribution is to track the data dependence in the residual terms
% $\resid{\phi}(\x_n)$ such that, when $\phi(\theta, \x_n)$ is third-order
% BCLT-okay, the residual is square summable with $\fdist$-probability approaching
% one.

% Specifically \thmref{bayes_clt_main} follows from \lemref{bayes_clt}, which
% establishes the expansion above with an explicit, non-stochastic  expression for
% the residual $\resid{\phi}(\x_n)$, and \lemref{resid_op1}, which shows that
% $\resid{\phi}(\x_n)$ is square-summable.
% %
% \end{proofsketch}



%%%%%%%%%%%%%%%%%%%%%%%%%%%%%%%%%%%%%%%%%
%%%%%%%%%%%%%%%%%%%%%%%%%%%%%%%%%%%%%%%%%
%%%%%%%%%%%%%%%%%%%%%%%%%%%%%%%%%%%%%%%%%




% First a classical FCLT
To connect \thmref{bayes_clt_main} to the IJ covariance, we first show that
\assuref{bayes_clt} is sufficient (indeed, stronger than necessary) to prove the
(known) results of \eqref{fclt_bclt} for third--order BCLT-okay functions. For completeness,
we state and prove these known results under  our stated assumptions.


\begin{assu}\assulabel{bclt_okay}
    The functions $\g(\theta)$, $\ell(\xn \vert \theta)$, and
    $\g(\theta) \ell(\xn \vert \theta)$ are all third-order BCLT-okay.
\end{assu}


%%%%%%%%%%%%%%%%%%%%%
%%%%%%%%%%%%%%%%%%%%%
\begin{prop}\proplabel{freq_clt}(See proof in \appref{freq_clt_proof},
    as well as \cite[Chapter 7]{van:2000:asymptotic})
%
Let \assuref{bayes_clt,bclt_okay} hold.
Then $\thetahat \plim \thetatrue$, and
%
\begin{align*}
\sqrt{N}\left(g(\thetahat) - g(\thetatrue)\right)
\dlim
\normdist\left(0, \gcovtrue \right)
\quad\textrm{where}\quad
\gcovtrue =
    \ggrad{1}(\thetatrue)
    \info^{-1} \scorecov \info^{-1} \ggrad{1}(\thetatrue)^\trans.
\end{align*}
%
\end{prop}
%
% \begin{proofsketch}{\appref{freq_clt_proof}}
% %
% \Propref{freq_clt} is only a minor variation of the classical proof of
% asymptotic normality of the maximum likelihood estimator, accounting for the
% possibility that, under misspecification, we may have $\info^{-1} \ne
% \scorecov$. \propref{freq_clt} is proven under weaker conditions than ours in,
% for example, \cite[Chapter 7]{van:2000:asymptotic}.  In \appref{freq_clt_proof}
% we include our own proof under our stated conditions for completeness.
% % conditions in the corollary to Theorem 3 of \citet{huber:1967:nonstandard}.
% %
% \end{proofsketch}%


% The limiting frequentist distribution of $\expect{\post}{g(\theta)}$ under
% misspecification, proven under weaker conditions than ours in \cite[Theorem
% 10.8]{van:2000:asymptotic}, follows from \thmref{bayes_clt_main}.  We state and
% prove the result here, since it establishes the estimand which the IJ covariance
% is targeting.

The next result shows that
% it follows from \thmref{bayes_clt_main} and
% \propref{freq_clt} that
our target variance is $\gcovtrue$.

%%%%%%%%%%%%%%%%%%%%%%%%%%%%%%%%%%%%%%%%%
%%%%%%%%%%%%%%%%%%%%%%%%%%%%%%%%%%%%%%%%%
\begin{cor}\corlabel{bayes_expectation_dist}(\cite[Theorem 10.8]{van:2000:asymptotic})
%
Under \assuref{bayes_clt,bclt_okay},
%
\begin{align*}
%
\sqrt{N} \left( \expect{\post}{g(\theta)} - g(\thetatrue) \right)
    \dlim \normdist\left(0, \gcovtrue\right).
%
\end{align*}
%
% where
% %
% \begin{align*}
% \gcovtrue := \ggrad{1}(\thetatrue) \info^{-1} \scorecov
%             \info^{-1} \ggrad{1}(\thetatrue)^T.
% \end{align*}
%
\begin{proof}
%
Under \assuref{bayes_clt,bclt_okay}, \lemref{regular_event} in \appref{bayes_clt_proof}
yields that
$\ggrad{1}(\thetahat)$, $\infohat$, and $\likhatk{3}(\thetahat)$ are all
$O_p(1)$.
%
By \thmref{bayes_clt_main} we thus have,
%
\begin{align*}
%
\sqrt{N} \left( \expect{\post}{g(\theta)} - g(\thetatrue) \right) =&
\sqrt{N} \left( g(\thetahat) - g(\thetatrue) \right) + O_p(N^{-1/2}),
%
\end{align*}
%
The distributional limit then follows from \propref{freq_clt} and Slutsky's
theorem.
%
\end{proof}
\end{cor}
%%%%%%%%%%%%%%%%%%%%%%%%%%%%%%%%%%%%%%%%%

Finally, our main result, the consistency of $\gcovij$, follows from
\thmref{bayes_clt_main} and \corref{bayes_expectation_dist}.

%%%%%%%%%%%%%%%%%%%%%%%%%%%%%%%%%%%%%
%%%%%%%%%%%%%%%%%%%%%%%%%%%%%%%%%%%%%
\begin{thm}\thmlabel{ij_consistent}(See proof in \appref{ij_consistent}.)
%
Under \assuref{bayes_clt,bclt_okay} the infinitesimal jackknife
variance estimate is consistent.  Specifically,
%
\begin{align*}
    \gcovij ={}&
        \ggrad{1}(\thetahat) \infohat^{-1}
            \scorecovhat \infohat^{-1} \ggrad{1}(\thetahat)^\trans +
        \ordlogp{N^{-1}}
        \plim \gcovtrue.
\end{align*}
%
\end{thm}
%%%%%%%%%%%%%%%%%%%%%%%%%%%%%%%%%%%%%
